% Options for packages loaded elsewhere
\PassOptionsToPackage{unicode}{hyperref}
\PassOptionsToPackage{hyphens}{url}
%
\documentclass[
  11pt,
  a4paper,
]{article}
\usepackage{amsmath,amssymb}
\usepackage{iftex}
\ifPDFTeX
  \usepackage[T1]{fontenc}
  \usepackage[utf8]{inputenc}
  \usepackage{textcomp} % provide euro and other symbols
\else % if luatex or xetex
  \usepackage{unicode-math} % this also loads fontspec
  \defaultfontfeatures{Scale=MatchLowercase}
  \defaultfontfeatures[\rmfamily]{Ligatures=TeX,Scale=1}
\fi
\usepackage{lmodern}
\ifPDFTeX\else
  % xetex/luatex font selection
  \setmainfont[]{DejaVu Serif}
  \setsansfont[]{DejaVu Sans}
  \setmonofont[]{DejaVu Sans Mono}
\fi
% Use upquote if available, for straight quotes in verbatim environments
\IfFileExists{upquote.sty}{\usepackage{upquote}}{}
\IfFileExists{microtype.sty}{% use microtype if available
  \usepackage[]{microtype}
  \UseMicrotypeSet[protrusion]{basicmath} % disable protrusion for tt fonts
}{}
\makeatletter
\@ifundefined{KOMAClassName}{% if non-KOMA class
  \IfFileExists{parskip.sty}{%
    \usepackage{parskip}
  }{% else
    \setlength{\parindent}{0pt}
    \setlength{\parskip}{6pt plus 2pt minus 1pt}}
}{% if KOMA class
  \KOMAoptions{parskip=half}}
\makeatother
\usepackage{xcolor}
\usepackage[margin=22mm]{geometry}
\usepackage{longtable,booktabs,array}
\usepackage{calc} % for calculating minipage widths
% Correct order of tables after \paragraph or \subparagraph
\usepackage{etoolbox}
\makeatletter
\patchcmd\longtable{\par}{\if@noskipsec\mbox{}\fi\par}{}{}
\makeatother
% Allow footnotes in longtable head/foot
\IfFileExists{footnotehyper.sty}{\usepackage{footnotehyper}}{\usepackage{footnote}}
\makesavenoteenv{longtable}
\setlength{\emergencystretch}{3em} % prevent overfull lines
\providecommand{\tightlist}{%
  \setlength{\itemsep}{0pt}\setlength{\parskip}{0pt}}
\setcounter{secnumdepth}{-\maxdimen} % remove section numbering
\ifLuaTeX
\usepackage[bidi=basic]{babel}
\else
\usepackage[bidi=default]{babel}
\fi
\babelprovide[main,import]{russian}
\ifPDFTeX
\else
\babelfont{rm}[]{DejaVu Serif}
\fi
% get rid of language-specific shorthands (see #6817):
\let\LanguageShortHands\languageshorthands
\def\languageshorthands#1{}
\usepackage{ragged2e}
\usepackage{needspace}
\usepackage{etoolbox}
\usepackage{xurl}
\hfuzz=0.3pt
\AtBeginDocument{%
  \RaggedRight
  \setlength{\emergencystretch}{3em}%
  \hypersetup{pdftitle={LRX: добавление нуля и неограниченный перерасход проекции}}%
}
\BeforeBeginEnvironment{longtable}{\Needspace{8\baselineskip}}
\AtBeginEnvironment{longtable}{\small}
\BeforeBeginEnvironment{itemize}{\Needspace{4\baselineskip}}
\ifLuaTeX
  \usepackage{selnolig}  % disable illegal ligatures
\fi
\IfFileExists{bookmark.sty}{\usepackage{bookmark}}{\usepackage{hyperref}}
\IfFileExists{xurl.sty}{\usepackage{xurl}}{} % add URL line breaks if available
\urlstyle{same}
\hypersetup{
  pdflang={ru-RU},
  hidelinks,
  pdfcreator={LaTeX via pandoc}}

\author{}
\date{}

\begin{document}

\hypertarget{ux434ux43eux431ux430ux432ux43bux435ux43dux438ux435-ux43dux443ux43bux44f-ux442ux43eux447ux43dux430ux44f-ux440ux435ux434ux443ux43aux446ux438ux44f-ux438-ux43dux435ux43eux433ux440ux430ux43dux438ux447ux435ux43dux43dux44bux439-ux43fux435ux440ux435ux440ux430ux441ux445ux43eux434-ux43fux440ux43eux435ux43aux446ux438ux438}{%
\section{Добавление нуля: точная редукция и неограниченный перерасход
проекции}\label{ux434ux43eux431ux430ux432ux43bux435ux43dux438ux435-ux43dux443ux43bux44f-ux442ux43eux447ux43dux430ux44f-ux440ux435ux434ux443ux43aux446ux438ux44f-ux438-ux43dux435ux43eux433ux440ux430ux43dux438ux447ux435ux43dux43dux44bux439-ux43fux435ux440ux435ux440ux430ux441ux445ux43eux434-ux43fux440ux43eux435ux43aux446ux438ux438}}

21 сентября 2026 года. Общая гипотеза покрытия остаётся открытой.

\hypertarget{ux440ux435ux437ux443ux43bux44cux442ux430ux442ux44b-ux438-ux438ux445-ux43eux442ux43dux43eux448ux435ux43dux438ux435-ux43a-ux433ux438ux43fux43eux442ux435ux437ux435}{%
\subsection{1. Результаты и их отношение к
гипотезе}\label{ux440ux435ux437ux443ux43bux44cux442ux430ux442ux44b-ux438-ux438ux445-ux43eux442ux43dux43eux448ux435ux43dux438ux435-ux43a-ux433ux438ux43fux43eux442ux435ux437ux435}}

Пусть \(n=m+r\), \(m\ge8\), \(r\ge2\), а канонический корень равен
\((1,\ldots,m,0^r)\). Целевой бюджет можно записать как

\[T_m(n)=\binom{m+1}{2}+(r-1)(m-2)
=\binom n2-\frac{(r-1)(r+4)}2.\]

Поэтому общая рекурсия

\[E_r(n)\le E_{r-1}(n-1)+m-2 \tag{1}\]

вместе с базовой оценкой \(E_1(m+1)\le\binom{m+1}{2}\) доказала бы всю
гипотезу. Поточечная версия через расстояние конкретного удалённого
вектора уже \href{LRX_MULTISET_ZERO_DELETION_AUDIT_RU.md}{опровергнута}.
Здесь исследуется другой подход: точный подъём путей меньшего графа.

Получены следующие результаты.

\begin{itemize}
\tightlist
\item
  Доказана точная динамика для путей с ограниченной длиной проекции
  после удаления одного помеченного нуля. Это работает для всех
  размеров.
\item
  Подъёма только кратчайших путей недостаточно уже при \(n=10,m=8\).
\item
  Ограниченные обходы в меньшем графе дают проверенные конечные
  сертификаты бюджетов \(42,48,54,52\) для четырёх наборов параметров.
  Эти численные границы ранее были известны из BFS; нова редукция.
\item
  Построено бесконечное семейство при \(r=2\), для которого минимально
  необходимый перерасход проекции равен \(2k-2\). Никакой постоянный
  перерасход не покрывает все размеры в этой модели.
\item
  Для того же семейства доказано точное расстояние \(6k-2\) и дано явное
  LRX-слово. Семейство удовлетворяет исходной гипотезе.
\end{itemize}

Рекурсия (1) остаётся открытой. Контрпример последнего раздела относится
к ограниченному классу проекций, а не к (1) или гипотезе покрытия.

\hypertarget{ux443ux434ux430ux43bux435ux43dux438ux435-ux43fux43eux43cux435ux447ux435ux43dux43dux43eux433ux43e-ux43dux443ux43bux44f}{%
\subsection{2. Удаление помеченного
нуля}\label{ux443ux434ux430ux43bux435ux43dux438ux435-ux43fux43eux43cux435ux447ux435ux43dux43dux43eux433ux43e-ux43dux443ux43bux44f}}

Временно отличим один ноль и обозначим его \(\bullet\). Состояние
задаётся парой \((u,j)\): \(u\) --- вектор длины \(n-1\), полученный
удалением \(\bullet\), а \(j\in\{0,\ldots,n-1\}\) --- его позиция
вставки. Пусть \(d(u)\) --- точное расстояние от \(u\) до меньшего корня
\(u_0=(1,\ldots,m,0^{r-1})\).

Для каждого шага переход пары имеет вид:

\begin{longtable}[]{@{}lll@{}}
\toprule\noalign{}
Шаг & Условие & Новая пара \\
\midrule\noalign{}
\endhead
\bottomrule\noalign{}
\endlastfoot
\(L\) & \(j=0\) & \((u,n-1)\) \\
\(L\) & \(j>0\) & \((Lu,j-1)\) \\
\(R\) & \(j=n-1\) & \((u,0)\) \\
\(R\) & \(j<n-1\) & \((Ru,j+1)\) \\
\(X\) & \(j=0\) или \(j=1\) & \((u,1-j)\) \\
\(X\) & \(j\ge2\) & \((Xu,j)\) \\
\end{longtable}

Это непосредственно следует из удаления помеченной позиции в исходном и
конечном списках. Если \(Xu=u\), последняя строка сохраняет и помеченное
состояние; такой шаг можно удалить.

\textbf{Лемма 1.} Удаление помеченного нуля переводит каждое ребро LRX в
ребро меньшего графа либо в неподвижный шаг. В частности, для каждого
вектора \(v\) и любой его нулевой позиции \(j\)

\[d(v\setminus j)\le d(v),\qquad E_{r-1}(n-1)\le E_r(n). \tag{2}\]

\textbf{Доказательство.} Следуем за выбранным нулём вдоль любого
сортирующего слова. В конце он находится в нулевом суффиксе, поэтому
удаление даёт меньший корень. После удаления неподвижных шагов таблица
даёт путь длины не больше исходной. Для второй оценки вставим ноль в
любой вектор меньшего графа и применим первую. \(\square\)

\hypertarget{ux442ux43eux447ux43dux430ux44f-ux434ux438ux43dux430ux43cux438ux43aux430-ux441-ux431ux44eux434ux436ux435ux442ux43eux43c-ux43fux440ux43eux435ux43aux446ux438ux438}{%
\subsection{3. Точная динамика с бюджетом
проекции}\label{ux442ux43eux447ux43dux430ux44f-ux434ux438ux43dux430ux43cux438ux43aux430-ux441-ux431ux44eux434ux436ux435ux442ux43eux43c-ux43fux440ux43eux435ux43aux446ux438ux438}}

Проекцией слова называем последовательность непустых шагов меньшего
графа из таблицы. Соседние обратные шаги проекции здесь \textbf{не
сокращаются}: они могут соответствовать необходимым движениям в полном
векторе.

Пусть \(F_e(u,j)\) --- минимальная длина полного сортирующего слова,
проекция которого содержит не больше \(d(u)+e\) шагов, где \(e\ge0\) ---
целое. При отсутствии такого слова полагаем \(F_e(u,j)=+\infty\).
Помеченный ноль в конце разрешён на любой позиции \(j\ge m\).

Для перехода с изменением проекции \(u\to u'\) положим

\[c(u,u')=1+d(u')-d(u)\in\{0,1,2\}. \tag{3}\]

Вдоль пути до меньшего корня сумма этих величин равна

\[\sum c=\text{длина проекции}-d(u). \tag{4}\]

Это телескопическая сумма разностей расстояний.

При фиксированном \(u\) единственные полезные переходы без изменения
проекции образуют трёхвершинный путь

\[1\ \overset{X}{\longleftrightarrow}\ 0\
\overset{L,R}{\longleftrightarrow}\ n-1.\]

Остальные позиции изолированы. Пусть \(\rho(j,k)\) --- расстояние в этом
графе, равное \(+\infty\) между разными компонентами. Для каждого
\((e,u,j)\) определим предварительную цену

\[A_e(u,j)=\min\left(
\{0:u=u_0,\ j\ge m\}\ \cup\
\{1+F_{e-c(u,u')}(u',j')\}\right), \tag{5}\]

где вторая часть минимума пробегает строки таблицы, меняющие \(u\), для
которых \(c(u,u')\le e\). Минимум пустого множества равен бесконечности.
Тогда точная рекурсия имеет вид

\[\boxed{F_e(u,j)=\min_k\bigl(\rho(j,k)+A_e(u,k)\bigr).} \tag{6}\]

\textbf{Теорема 1.} Формулы (3)--(6) вычисляют именно указанный минимум
по всем допустимым словам, а не эвристическую верхнюю оценку.

\textbf{Доказательство.} До первого изменения проекции слово движется
только в графе \(\rho\). Затем либо достигается корень, либо выполняется
один переход из (5). Оставшийся бюджет равен \(e-c\) по (4). Обратно,
каждый член (5) и путь в \(\rho\) реализуются буквами LRX из таблицы.
Поэтому оба направления неравенства в (6) выполнены.

Рекурсия корректно упорядочена по паре \((e,d(u))\). Если \(c=0\), то
\(d(u')=d(u)-1\); если \(c>0\), уменьшается \(e\). Переходы внутри
\(\rho\) обработаны точным конечным минимумом. Таким образом,
циклических зависимостей вычисления нет. \(\square\)

Для видимого вектора положим

\[B_e(v)=\min_{j:v_j=0}F_e(v\setminus j,j).\]

Всегда \(d(v)\le B_e(v)\). Более того,

\[e\ge d(v)-\max_{j:v_j=0}d(v\setminus j)
\quad\Longrightarrow\quad B_e(v)=d(v). \tag{7}\]

Действительно, выберем удаление, достигающее максимума, и проследим его
вдоль кратчайшего полного слова. Длина проекции не превосходит длины
этого слова, поэтому его перерасход укладывается в (7). Это свойство
объясняет полноту модели при растущем бюджете; само по себе оно не даёт
численной оценки радиуса.

\hypertarget{ux43aux43eux43dux435ux447ux43dux430ux44f-ux43fux440ux43eux432ux435ux440ux43aux430-ux440ux435ux43aux443ux440ux441ux438ux432ux43dux43eux439-ux43aux43eux43dux441ux442ux440ux443ux43aux446ux438ux438}{%
\subsection{4. Конечная проверка рекурсивной
конструкции}\label{ux43aux43eux43dux435ux447ux43dux430ux44f-ux43fux440ux43eux432ux435ux440ux43aux430-ux440ux435ux43aux443ux440ux441ux438ux432ux43dux43eux439-ux43aux43eux43dux441ux442ux440ux443ux43aux446ux438ux438}}

\href{scripts/multiset_marked_zero_geodesic.cpp}{Проверяющий C++}
использует BFS только в меньшем графе и затем вычисляет (6). В большем
графе он перечисляет векторы, но не вычисляет настоящий BFS.

Независимый \href{scripts/audit_multiset_marked_zero.py}{сценарий на
кортежах} строит обратный BFS по буквальным помеченным векторам с
бюджетом. Он не использует ранжирование или рекурсию C++. Все
гистограммы совпали на 50 проверках: \(4\le n\le7\), \(2\le m\le n-2\),
\(0\le e\le4\).

В следующей таблице «нет пути» означает отсутствие пути именно в
ограниченном классе, а не недостижимость в графе LRX. «Сверх бюджета»
учитывает только конечные значения \(B_e\).

\begin{longtable}[]{@{}
  >{\raggedright\arraybackslash}p{(\columnwidth - 8\tabcolsep) * \real{0.1579}}
  >{\raggedleft\arraybackslash}p{(\columnwidth - 8\tabcolsep) * \real{0.2105}}
  >{\raggedleft\arraybackslash}p{(\columnwidth - 8\tabcolsep) * \real{0.2105}}
  >{\raggedleft\arraybackslash}p{(\columnwidth - 8\tabcolsep) * \real{0.2105}}
  >{\raggedleft\arraybackslash}p{(\columnwidth - 8\tabcolsep) * \real{0.2105}}@{}}
\toprule\noalign{}
\begin{minipage}[b]{\linewidth}\raggedright
\(n,m,r\)
\end{minipage} & \begin{minipage}[b]{\linewidth}\raggedleft
\(e\)
\end{minipage} & \begin{minipage}[b]{\linewidth}\raggedleft
Нет пути
\end{minipage} & \begin{minipage}[b]{\linewidth}\raggedleft
Превышений бюджета
\end{minipage} & \begin{minipage}[b]{\linewidth}\raggedleft
Максимум конечного \(B_e\)
\end{minipage} \\
\midrule\noalign{}
\endhead
\bottomrule\noalign{}
\endlastfoot
\(10,8,2\) & 0 & 1524 & 83 & 46 \\
\(10,8,2\) & 1 & 439 & 0 & 42 \\
\(10,8,2\) & 2 & 12 & 0 & 42 \\
\(10,8,2\) & 3 & 0 & 0 & 42 \\
\(11,8,3\) & 0 & 0 & 33 & 50 \\
\(11,8,3\) & 2 & 0 & 0 & 48 \\
\(12,8,4\) & 0 & 0 & 95 & 57 \\
\(12,8,4\) & 2 & 0 & 0 & 54 \\
\(11,9,2\) & 0 & 5000 & 73 & 56 \\
\(11,9,2\) & 2 & 14 & 0 & 52 \\
\(11,9,2\) & 4 & 0 & 0 & 52 \\
\end{longtable}

Большие проверки охватили 48384000 различных видимых векторов в четырёх
графах. Независимо исполнены 2412 выданных контрольных слов, суммарно
27774 буквы; для каждого проверена и длина проекции. Полный
\href{literature/multiset_marked_zero_lift_final_20260920.json}{отчёт}
содержит гистограммы, слова и хеши исходников. Проверка начата 20
сентября и завершена 21 сентября; имя отчёта сохранено.

Пример ограничения геодезического подъёма:

\[v=(2,5,1,0,6,7,0,4,3,8).\]

Два удаления имеют расстояния \(24\) и \(28\) в меньшем графе. При
\(e=0\) первый подъём отсутствует, второй стоит \(46\). При \(e=3\)
имеется слово длины \(33\):

\begin{verbatim}
LXLLXLXLLLLXLXRRXRXRXLLLXLXLXRXRX
\end{verbatim}

Для удаления первой нулевой позиции его проекция имеет длину \(27\), для
второй --- \(29\). Настоящее расстояние здесь не объявляется равным
\(33\); для проверки гипотетического бюджета \(42\) достаточно верхнего
слова.

\hypertarget{ux431ux435ux441ux43aux43eux43dux435ux447ux43dux43eux435-ux43fux440ux435ux43fux44fux442ux441ux442ux432ux438ux435-ux434ux43bux44f-ux43fux43eux441ux442ux43eux44fux43dux43dux43eux433ux43e-ux43fux435ux440ux435ux440ux430ux441ux445ux43eux434ux430}{%
\subsection{5. Бесконечное препятствие для постоянного
перерасхода}\label{ux431ux435ux441ux43aux43eux43dux435ux447ux43dux43eux435-ux43fux440ux435ux43fux44fux442ux441ux442ux432ux438ux435-ux434ux43bux44f-ux43fux43eux441ux442ux43eux44fux43dux43dux43eux433ux43e-ux43fux435ux440ux435ux440ux430ux441ux445ux43eux434ux430}}

\textbf{Теорема 2.} Пусть \(k\ge2\), \(m\ge8k-1\), \(n=m+2\), \(r=2\), и

\[v_{m,k}=(1,\ldots,k,0,k+1,\ldots,m-k+1,0,m-k+2,\ldots,m). \tag{8}\]

Тогда:

\begin{enumerate}
\def\labelenumi{\arabic{enumi}.}
\tightlist
\item
  удаление любого из двух нулей даёт расстояние \(3k-1\);
\item
  \(B_e(v_{m,k})=+\infty\) при \(e<2k-2\);
\item
  при \(e=2k-2\) имеется полное сортирующее слово длины \(6k-2\);
\item
  настоящее LRX-расстояние \(d(v_{m,k})\) равно \(6k-2\).
\end{enumerate}

Таким образом, минимальный перерасход проекции, необходимый хотя бы для
существования подъёма, равен \(2k-2\) и не ограничен константой. Вся
зависимость от бюджета имеет вид

\[B_e(v_{m,k})=
\begin{cases}
+\infty,&0\le e<2k-2,\\
6k-2,&e\ge2k-2.
\end{cases}\]

\hypertarget{ux43bux43eux43aux430ux43bux438ux437ux430ux446ux438ux44f-ux43aux43eux440ux43eux442ux43aux43eux433ux43e-ux441ux43bux43eux432ux430}{%
\subsubsection{5.1. Локализация короткого
слова}\label{ux43bux43eux43aux430ux43bux438ux437ux430ux446ux438ux44f-ux43aux43eux440ux43eux442ux43aux43eux433ux43e-ux441ux43bux43eux432ux430}}

Используем неподвижную физическую окружность. Курсор начинает в \(0\);
\(L,R\) двигают его на \(+1,-1\), \(X\) меняет токены в курсоре и справа
от него. Если слово содержит \(R_*\) вращений, то затронуто обменами не
больше \(R_*+2\) физических позиций: курсор посещает связную дугу не
более чем из \(R_*+1\) позиций, а обмены добавляют только её правую
соседнюю позицию.

Следовательно, среди более чем \(\ell+2\) токенов с одинаковым исходным
смещением хотя бы один остаётся нетронутым в слове длины \(\ell\). Он
определяет конечную физическую фазу: если токен \(x\) исходно стоит в
\(x-1+c\) и не перемещён, то конечный курсор обязан быть \(c\) по модулю
длины окружности.

Во всех применениях ниже длина слова меньше длины окружности. Указанные
позиции около \(0\) имеют единственные доступные целые представители.
Поэтому пути курсора оцениваются как пути на прямой.

\hypertarget{ux434ux432ux430-ux443ux434ux430ux43bux451ux43dux43dux44bux445-ux432ux435ux43aux442ux43eux440ux430}{%
\subsubsection{5.2. Два удалённых
вектора}\label{ux434ux432ux430-ux443ux434ux430ux43bux451ux43dux43dux44bux445-ux432ux435ux43aux442ux43eux440ux430}}

После удаления первого нуля получаем

\[u_R=(1,\ldots,m-k+1,0,m-k+2,\ldots,m).\]

На окружности длины \(M=m+1\) ноль стоит в \(-k\). Все \(m-k+1\)
начальных меток стоят в своих физических целях. Для слова длины не
больше \(5k-4\) число таких меток больше \(\ell+2\), поэтому конечный
курсор равен \(0\).

Ноль должен попасть в \(-1\). Обход длинной стороны потребовал бы больше
\(5k-4\) обменов. Значит, нужно хотя бы \(k-1\) обменов, а курсор обязан
посетить \(-k\) и вернуться в \(0\): хотя бы \(2k\) вращений. Нижняя
оценка равна \(3k-1\) и достигается словом

\[R^k X(LX)^{k-2}L^2. \tag{9}\]

После удаления второго нуля получаем

\[u_L=(1,\ldots,k,0,k+1,\ldots,m).\]

Здесь \(m-k\) меток после нуля стоят в физических позициях \(x\), а
должны оказаться в \(x-1+c\). Локализация короткого слова вынуждает
\(c=1\). Ноль идёт из \(k\) в \(0\), поэтому нужны хотя бы \(k\)
обменов. Курсор посещает \(k-1\), затем \(0\) и заканчивает в \(1\):
вращений хотя бы \(2k-1\). Оценка \(3k-1\) достигается словом

\[L^{k-1}X(RX)^{k-1}L. \tag{10}\]

Эти нижние оценки исключают все более короткие слова и доказывают первый
пункт теоремы.

\hypertarget{ux43eux431ux44fux437ux430ux442ux435ux43bux44cux43dux43eux435-ux43fux43eux441ux435ux449ux435ux43dux438ux435-ux443ux434ux430ux43bux451ux43dux43dux43eux433ux43e-ux43fux440ux43eux43cux435ux436ux443ux442ux43aux430}{%
\subsubsection{5.3. Обязательное посещение удалённого
промежутка}\label{ux43eux431ux44fux437ux430ux442ux435ux43bux44cux43dux43eux435-ux43fux43eux441ux435ux449ux435ux43dux438ux435-ux443ux434ux430ux43bux451ux43dux43dux43eux433ux43e-ux43fux440ux43eux43cux435ux436ux443ux442ux43aux430}}

Пусть существует подъём с \(e<2k-2\). Его проекция имеет длину

\[\ell\le(3k-1)+(2k-3)=5k-4. \tag{11}\]

До первого попадания пометки в позиции \(0,1,n-1\) каждый шаг проекции
совпадает с полным шагом. Если такого попадания вообще нет, то
помеченный ноль не меняет физического положения обменами. Конечная фаза
проекции из §5.2 равна \(0\) либо \(1\), и пометка остаётся среди
различных меток, а не в нулевом суффиксе. Это невозможно.

Для первого нуля исходная позиция вставки равна \(j=k\). Первое
посещение управляющей тройки требует от курсора проекции дойти до
\(j-1=k-1\) либо до \(j-M\). Второй вариант дальше бюджета (11).
Следовательно, при сортировке \(u_R\) курсор должен посетить и \(-k\), и
\(k-1\), начиная и заканчивая в \(0\). Это требует хотя бы

\[2\bigl(k+(k-1)\bigr)=4k-2\]

вращений и хотя бы \(k-1\) обменов. Длина проекции не меньше \(5k-3\), в
противоречии с (11).

Для второго нуля позиция вставки равна \(j=n-k=M+1-k\). Теперь первое
достижимое управляющее положение требует посещения \(j-M=-(k-1)\)
курсором проекции. При сортировке \(u_L\) ноль идёт влево из \(k\) в
\(0\), поэтому курсор обязан посетить \(k-1\) \textbf{до} посещения
\(0\) на соответствующих обменах. Конечный курсор равен \(1\). Посещение
\(-(k-1)\) может быть до этой пары событий, между ними или после. В
каждом случае нижняя оценка числа вращений равна \(4k-3\); например,

\[0\longrightarrow -(k-1)\longrightarrow k-1
\longrightarrow0\longrightarrow1.\]

С хотя бы \(k\) обменами снова получаем \(5k-3>\ell\). Таким образом,
оба возможных удаления исключены. Это доказывает пункт 2.

\hypertarget{ux44fux432ux43dux43eux435-ux441ux43bux43eux432ux43e-ux438-ux442ux43eux447ux43dux43eux435-ux43fux43eux43bux43dux43eux435-ux440ux430ux441ux441ux442ux43eux44fux43dux438ux435}{%
\subsubsection{5.4. Явное слово и точное полное
расстояние}\label{ux44fux432ux43dux43eux435-ux441ux43bux43eux432ux43e-ux438-ux442ux43eux447ux43dux43eux435-ux43fux43eux43bux43dux43eux435-ux440ux430ux441ux441ux442ux43eux44fux43dux438ux435}}

Рассмотрим слово

\[\boxed{W_k=L^{k-1}X(RX)^{k-1}R^kX(LX)^{k-2}L^3.} \tag{12}\]

Первый блок переносит ноль из \(k\) в \(0\) и заканчивает курсором в
\(0\). Затем \(R^k\) подводит курсор ко второму нулю в \(-k\). Второй
блок переносит его в \(-1\) и заканчивает курсором в \(-2\). Последние
три \(L\) устанавливают курсор в \(1\). Метки \(1,\ldots,m\) занимают
физические позиции \(1,\ldots,m\), нули --- \(-1,0\). Видимый вектор
равен каноническому корню.

В слове \(2k-1\) обменов и \(4k-1\) вращений, всего \(6k-2\) букв. При
удалении любого из двух помеченных нулей исчезает \(k+1\) шагов, так что
длина проекции равна \(5k-3\). Её перерасход равен \(2k-2\). Это
доказывает пункт 3.

Осталось доказать нижнюю границу полного расстояния. Допустим, имеется
сортирующее слово длины \(\ell\le6k-3\). Среди средних меток
\(k+1,\ldots,m-k+1\) есть \(m-2k+1\ge6k\) токенов с исходной физической
позицией \(x\). Так как \(\ell+2\le6k-1\), один из них не перемещён.
Конечный курсор поэтому равен \(1\).

Целевые нулевые позиции равны \(-1,0\). Ноль из \(k\) должен идти влево,
а ноль из \(-k\) --- вправо: длинные обходы превышают весь бюджет.
Обмены двух нулей можно удалить, поскольку они не меняют видимый вектор.
Поэтому число обменов не меньше \(k+(k-1)=2k-1\).

Курсор обязан посетить \(k-1\), затем \(0\) при переносе первого нуля, а
также \(-k\) при переносе второго. Он начинает в \(0\) и заканчивает в
\(1\). Независимо от положения посещения \(-k\) относительно
упорядоченной пары \(k-1,0\), путь требует хотя бы \(4k-1\) вращений.
Например,

\[0\longrightarrow-k\longrightarrow k-1\longrightarrow0\longrightarrow1\]

имеет именно такую длину. Всего нужно хотя бы \(6k-2\) шагов, что
противоречит выбору \(\ell\). Пункт 4 и теорема доказаны.

Семейство (8) лежит в области исходной гипотезы и удовлетворяет ей:

\[d(v_{m,k})=6k-2\le\binom{m+2}{2}-3=T_m(m+2).\]

Таким образом, фиксированный перерасход проекции недостаточен даже для
некоторых векторов, расстояние которых значительно ниже бюджета.

\hypertarget{ux432ux43eux441ux43fux440ux43eux438ux437ux432ux43eux434ux438ux43cux43eux441ux442ux44c-ux438-ux441ux43bux435ux434ux443ux44eux449ux438ux439-ux43eux442ux43aux440ux44bux442ux44bux439-ux448ux430ux433}{%
\subsection{6. Воспроизводимость и следующий открытый
шаг}\label{ux432ux43eux441ux43fux440ux43eux438ux437ux432ux43eux434ux438ux43cux43eux441ux442ux44c-ux438-ux441ux43bux435ux434ux443ux44eux449ux438ux439-ux43eux442ux43aux440ux44bux442ux44bux439-ux448ux430ux433}}

\href{scripts/audit_multiset_marked_zero_barrier.py}{Независимый аудит
семейства} проверил 297 входов при \(2\le k\le100\) и трёх значениях
\(m\) для каждого \(k\). Исполнены 594 помеченных слова, суммарно 180576
букв.

Отдельная разреженная динамика по длине проекции, без C++-ранжировщика,
подтвердила отсутствие подъёмов при \(k=2,3,4\), \(m=8k-1\), с бюджетами
проекции \(6,11,16\). Полные шары меньшего графа имели 185, 4405 и 85610
вершин соответственно.

Точные полные расстояния независимо проверены словами и непересечением
двух полных шаров:

\begin{longtable}[]{@{}rrrrll@{}}
\toprule\noalign{}
\(k\) & \(m\) & \(n\) & Расстояние & Радиусы шаров & Размеры шаров \\
\midrule\noalign{}
\endhead
\bottomrule\noalign{}
\endlastfoot
2 & 15 & 17 & 10 & \(4,5\) & \(42,94\) \\
3 & 23 & 25 & 16 & \(7,8\) & \(307,682\) \\
4 & 31 & 33 & 22 & \(10,11\) & \(1989,4405\) \\
5 & 39 & 41 & 28 & \(13,14\) & \(11990,26536\) \\
\end{longtable}

В каждой строке сумма радиусов на единицу меньше расстояния.
\href{literature/multiset_marked_zero_barrier_20260921.json}{Протокол}.
Перенос на все \(k,m\) следует из текстового доказательства §5.

Команды проверки из корня проекта:

\begin{verbatim}
g++ -std=c++17 -O3 -Wall -Wextra \
  scripts/multiset_marked_zero_geodesic.cpp -o /tmp/lrx_marked_zero_check
python scripts/audit_multiset_marked_zero.py \
  --binary /tmp/lrx_marked_zero_check --large \
  --output /tmp/lrx_marked_zero_new_audit.json
python scripts/audit_multiset_marked_zero_barrier.py \
  --output /tmp/lrx_marked_zero_new_barrier.json
\end{verbatim}

Предел C++-таблицы явно ограничен 500 миллионами двухбайтовых значений;
дополнительная память нужна для меньшего BFS. Кортежные шары ограничены
300000 вершинами. Выходные файлы создаются только под новыми именами.

Общая рекурсия (1) пока не доказана. Проверки показывают, что запас
между \(d(u)\) и общим радиусом \(E_{r-1}(n-1)\) может оплачивать
длинные обходы проекции. Теорема 2 исключает замену такого зависящего от
состояния бюджета одной универсальной константой. Для завершения
гипотезы нужен общий контроль полной цены подъёма с этим запасом либо
другая конструкция.

Один точный следующий кандидат получается при \(P=E_{r-1}(n-1)\):

\[A_P(v)=\min_{j:v_j=0}F_{P-d(v\setminus j)}(v\setminus j,j).\]

Индекс бюджета неотрицателен по определению \(P\). Это минимум среди
слов с длиной проекции не больше общего радиуса меньшего графа.
Утверждение \(A_P(v)\le P+m-2\) для всех \(v\) дало бы (1). Оно
допускает большой перерасход для удалённых векторов, близких к меньшему
корню, и пока остаётся открытым.

\end{document}
